% =========================================================================
% PLATFORM TAXONOMY ROUTING (patho.dentnet.space)
% =========================================================================
% DOC_ID: benign_epithelial
% DOC_SLUG: benign_epithelial
% DOC_CATEGORY: cat_benign_epithelial
% DOC_TITLE: الگوریتم جامع: ضایعات اپی‌تلیال خوش‌خیم و لکوپلاکیا
% DOC_ALIASES: benign epithelial, لکوپلاکیا, ضایعات اپی تلیال, پاپیلوما, کندیلوما, دیسپلازی اپی تلیال, leukoplakia, papilloma
% =========================================================================
% === CLINICAL EVIDENCE LEDGER & PMIDS ===
% 1. WHO Classification of Head and Neck Tumours (5th Ed, 2024) - Oral Epithelial Dysplasia & HPV-related Lesions.
% 2. Neville BW, Damm DD, Allen CM, Chi AC. Oral and Maxillofacial Pathology (5th Ed).
% 3. PMID: 42690553 - Khurram SA et al. (Head Neck Pathol 2026) Pattern-Anchored Diagnosis of Oral Epithelial Dysplasia: COST Action INTERCEPTOR.
% 4. PMID: 42969171 - Grzech-Leśniak K et al. (Dent Med Probl 2026) Laser-Based Management of Oral Leukoplakia: WFLD Study Group Consensus.
% 5. PMID: 42808717 - Zhao J et al. (Oral Health Prev Dent 2026) Clinical Efficacy of 5-Aminolevulinic Acid Photodynamic Therapy in Oral Leukoplakia.
% 6. PMID: 41050292 - Goswami M et al. (Int J Clin Pediatr Dent 2025) Minimally Invasive Management of Oral Squamous Papilloma.
% 7. PMID: 40597151 - Huang J et al. (BMC Oral Health 2025) Recurrent Verrucous Carcinoma and Proliferative Verrucous Leukoplakia Risk.
% 8. PMID: 39290142 - Bhattarai BP et al. (J Dent Res 2024) Recurrence in Oral Leukoplakia: Systematic Review and Meta-analysis.
% =========================================================================

\documentclass[a4paper]{article}
\usepackage[margin=1.15cm,footskip=0pt,headsep=0pt]{geometry}
\usepackage{amsmath,amssymb}
\usepackage{xcolor}
\usepackage{tikz}
\usetikzlibrary{arrows.meta,positioning,shapes.geometric}
\usepackage{hyperref}
\usepackage{fontspec}

\makeatletter
\newcommand{\MathCollectFalse}{}
\newcommand{\MathCollectTrue}{}
\let\@kernel@tabular@init\relax
\makeatother

\usepackage{xepersian}

\settextfont[
  Path = Peyda/,
  BoldFont = Peyda-Bold.ttf,
  ItalicFont = Peyda-Medium.ttf,
  BoldItalicFont = Peyda-Bold.ttf
]{Peyda-Regular.ttf}

\setlatintextfont{Times New Roman}

\IfFileExists{Peyda/Symbola.ttf}{%
  \newfontfamily\symb[Path=Peyda/]{Symbola.ttf}%
}{%
  \IfFileExists{editor/Peyda/Symbola.ttf}{%
    \newfontfamily\symb[Path=editor/Peyda/]{Symbola.ttf}%
  }{%
    \newfontfamily\symb{Symbola}%
  }%
}
\newcommand{\emoji}[1]{{\symb #1}}

\linespread{1.10}
\pagestyle{empty}
\thispagestyle{empty}

\raggedbottom
\setlength{\parindent}{0pt}
\setlength{\parskip}{0pt}

\definecolor{hlgray}{HTML}{E2E8F0}
\definecolor{hlyellow}{HTML}{FEF08A}
\definecolor{hlred}{HTML}{FEE2E2}
\definecolor{darkred}{HTML}{991B1B}
\definecolor{navy}{HTML}{0A2540}
\definecolor{teal}{HTML}{0D9488}

\begin{document}

\begin{center}
{\Large\bfseries فلوچارت تصمیم‌گیری بالینی: ضایعات اپی‌تلیال خوش‌خیم و لکوپلاکیا}\\[2pt]
{\normalsize\bfseries\color{navy} \lr{Benign Epithelial Lesions and Oral Leukoplakia (Differential Triage)}}\\[3pt]
{\small\itshape تریاژ پاپیلوماهای ویروسی (\lr{HPV})، لکوپلاکیای دیسپلاستیک (\lr{WHO 5th Ed})، مارجین اکسیزیون لیزر و پایش بدخیمی}
\end{center}
\vspace{-4pt}
\hrule height 0.8pt
\vspace{6pt}

\begin{center}
\begin{latin}
\linespread{0.95}\selectfont
\begin{tikzpicture}[
    every node/.style={font=\sffamily\footnotesize, align=center},
    stepbox/.style={rectangle, draw=black, fill=white, line width=0.7pt, inner sep=4.5pt},
    decision/.style={diamond, aspect=2.4, draw=black, fill=white, line width=0.7pt, inner sep=2pt, font=\sffamily\scriptsize\bfseries},
    arrow/.style={-{Stealth[scale=0.9]}, thick, black}
]

% 1. Initial Presentation
\node (n1) [stepbox, line width=1.0pt, text width=13.0cm] at (0, 0.0) {
    \rl{\textbf{تظاهر بالینی: پلاک سفید بدون قابلیت برداشت، برجستگی پاپیلاری، اگزوفیتیک یا پدانکوله}}\\[2pt]
    \scriptsize \rl{برجستگی گل‌کلمی مخاطی $\bullet$ پلاک کراتوتیک کف دهان، زبان یا لثه $\bullet$ لکوپلاکیای هموژن یا غیرهموژن $\bullet$ رد تروما}
};

% 2. Mandatory Step 1
\node (n2) [stepbox, text width=14.2cm] at (0, -1.60) {
    \rl{\colorbox{hlyellow}{\strut\textbf{مرحله ۱ الزامی: رد عوامل اتولوژی واکنشی (اصطکاک/دندان تیز) و ارزیابی فتوگرافی دیجیتال}}}\\[2pt]
    \scriptsize \rl{برطرف کردن لبه تیز دندان یا دنچر به مدت ۲ هفته؛ در صورت پایدار ماندن پلاک کراتوتیک $\rightarrow$ نیاز قطعی به ارزیابی بیوپسی}
};

% 3. Solitary Pedunculated Decision
\node (d1) [decision] at (0, -3.30) {\rl{ضایعه منفرد، پدانکوله،}\\[-1pt]\scriptsize \lr{Pedunculated / Warty?}};
\node (r_benign) [stepbox, draw=teal, line width=0.9pt, text width=3.8cm] at (5.8, -3.30) {
    \rl{\textbf{پاپیلوما منفرد (\lr{HPV 6/11}) \emoji{🔍}}}\\[2pt]
    \scriptsize \rl{ضایعه خوش‌خیم بدون دیسپلازی؛}\\\rl{اکسیزیون ساده با شالوده پدانکول؛}\\\colorbox{hlgray}{\strut\rl{\textbf{پیش‌آگهی عالی؛ عود نادر}}}
};

% 4. Mandatory Biopsy & Dysplasia Grading
\node (n3) [stepbox, line width=1.0pt, text width=15.8cm] at (0, -5.00) {
    \rl{\colorbox{hlyellow}{\strut\textbf{بیوپسی برشی (\lr{Incisional Biopsy}) در ضایعات وسیع یا نامنظم و گریدینگ دیسپلازی \emoji{🔬}}}}\\[2pt]
    \scriptsize \rl{نمونه‌برداری از بدنماترین ناحیه در اریترولکوپلاکیا $\bullet$ گریدینگ دیسپلازی اپی‌تلیال بر اساس معیارهای \lr{WHO 5th Ed 2024} \lr{[\textbf{PMID: 42690553}]}}
};

% 5. Subtype Decision
\node (d2) [decision] at (0, -6.60) {\rl{پاتولوژی و دیسپلازی}\\[-1pt]\scriptsize \lr{Dysplasia Grade?}};

% 6. 4 Main Columns
\node (h1) [stepbox, text width=3.6cm] at (-6.15, -9.00) {
    \rl{\textbf{پاپیلوما اسکواموس}}\\[2pt]
    \scriptsize \lr{Squamous Papilloma}\\\rl{پدانکوله با زوائد انگشتی؛}\\\rl{فاقد اتیپی یا دیسپلازی اپی‌تلیال؛}\\\rl{عود با جراحی: \colorbox{hlgray}{\strut\textbf{\rl{کمتر از ۱٪}}}}
};
\node (h2) [stepbox, text width=3.6cm] at (-2.05, -9.00) {
    \rl{\textbf{کوندیلوما / وروکا}}\\[2pt]
    \scriptsize \lr{Condyloma / Verruca}\\\rl{ضایعات سسیل پهن، مولتیپل؛}\\\rl{انتقال جنسی یا خودتلقیحی (\lr{HPV})؛}\\\rl{عود موضعی: \textbf{۱۰٪ تا ۲۰٪}}
};
\node (h3) [stepbox, text width=3.6cm] at (2.05, -9.00) {
    \rl{\textbf{دیسپلازی خفیف}}\\[2pt]
    \scriptsize \lr{Mild Dysplasia (WHO Low)}\\\rl{دیسپلازی محدود به ۱/۳ قاعده اپی‌تلیوم؛}\\\rl{پلاک هموژن با هایپرکراتوز پایدار؛}\\\rl{ریسک بدخیمی: \textbf{۳٪ تا ۵٪}}
};
\node (h4) [stepbox, draw=darkred, line width=1.2pt, text width=3.6cm] at (6.15, -9.00) {
    \rl{\textbf{دیسپلازی شدید / \lr{CIS} \emoji{⚡}}}\\[2pt]
    \scriptsize \lr{High-Grade Dysplasia / CIS}\\\rl{درگیری بیش از ۲/۳ تا تمام ضخامت؛}\\\rl{میتوز آتیپیک و پلئومورفیسم سلولی؛}\\\rl{ریسک بدخیمی:} \colorbox{hlred}{\strut\color{darkred}\textbf{\rl{۱۵٪ تا ۳۵٪}}}
};

% 7. Primary Interventions
\node (i1) [stepbox, text width=3.6cm] at (-6.15, -11.60) {
    \rl{\textbf{اکسیزیون کانوکشنال}}\\[2pt]
    \scriptsize \rl{برداشت جراحی با حاشیه سالم؛}\\\rl{شامل قاعده ساقه فیبروواسکولار؛}\\\rl{تایید پاتولوژی مارجین}\\\rl{\lr{[\textbf{PMID: 41050292}]}}
};
\node (i2) [stepbox, text width=3.6cm] at (-2.05, -11.60) {
    \rl{\textbf{اکسیزیون کامل $\pm$ کرایو}}\\[2pt]
    \scriptsize \rl{اکسیزیون با بیستوری یا لیزر \lr{CO2}؛}\\\rl{بررسی همزمان پارتنر جنسی؛}\\\rl{استفاده از مکش دود جراحی}\\\rl{(پیشگیری از انتقال ویروس با دود)}
};
\node (i3) [stepbox, text width=3.6cm] at (2.05, -11.60) {
    \rl{\textbf{تبخیر لیزر یا اکسیزیون}}\\[2pt]
    \scriptsize \rl{اکسیزیون با مارجین ۳ میلی‌متر، یا}\\\rl{تبخیر سطحی با لیزر \lr{CO2}؛}\\\rl{حذف کامل دخانیات و الکل}\\\rl{\lr{[\textbf{PMID: 42969171}]}}
};
\node (i4) [stepbox, draw=darkred, line width=1.2pt, text width=3.6cm] at (6.15, -11.60) {
    \rl{\textbf{رزکسیون با مارجین ۵mm}}\\[2pt]
    \scriptsize \textbf{\rl{اکسیزیون وسیع با حاشیه پاک ۵mm؛}}\\\rl{ارزیابی کامل مارجین بافتی؛}\\\rl{فتودینامیک‌تراپی ادجوانت (\lr{PDT})}\\\rl{\lr{[\textbf{PMID: 39290142}]}}
};

% 8. Surveillance Regimens
\node (s1) [stepbox, text width=3.6cm] at (-6.15, -14.20) {
    \rl{\textbf{زمان‌بندی پیگیری:}}\\[2pt]
    \rl{\colorbox{hlgray}{\strut\textbf{سالانه در معاینات روتین}}}\\[2pt]
    \scriptsize \rl{پیش‌آگهی عالی؛}\\\rl{عدم نیاز به تصویربرداری دوره‌ای}
};
\node (s2) [stepbox, text width=3.6cm] at (-2.05, -14.20) {
    \rl{\textbf{زمان‌بندی پیگیری:}}\\[2pt]
    \rl{\textbf{هر ۶ ماه تا ۲ سال}}\\[2pt]
    \scriptsize \rl{پایش دقیق پیدایش}\\\rl{ضایعات اگزوفیتیک و پاپیلاری جدید}
};
\node (s3) [stepbox, text width=3.6cm] at (2.05, -14.20) {
    \rl{\textbf{زمان‌بندی پیگیری:}}\\[2pt]
    \rl{\textbf{هر ۶ ماه (مادام‌العمر)}}\\[2pt]
    \scriptsize \rl{معاینه کامل مخاط و عکاسی دیجیتال؛}\\\rl{بیوپسی مجدد در صورت تغییر بالینی}
};
\node (s4) [stepbox, draw=darkred, line width=1.2pt, text width=3.6cm] at (6.15, -14.20) {
    \rl{\textbf{زمان‌بندی پیگیری:}}\\[2pt]
    \colorbox{hlred}{\strut\color{darkred}\textbf{\rl{هر ۳ ماه (۲ سال اول)}}}\\[2pt]
    \scriptsize \rl{سپس هر ۶ ماه مادام‌العمر؛}\\\rl{غربالگری دقیق تحول به \lr{OSCC}}
};

% 9. Clinical Red Flags
\node (redflags) [stepbox, draw=darkred, line width=0.9pt, text width=15.8cm] at (0, -17.40) {
    \rl{\textbf{هشدارهای بالینی و اندیکاسیون‌های ارجاع فوری یا بیوپسی مجدد (\lr{Clinical Red Flags}):}}\\[3pt]
    \scriptsize \colorbox{hlred}{\strut\color{darkred}\textbf{\rl{مداخله فوری تشخیصی در:}}} \rl{(۱) پیدایش پلاک قرمز (\lr{Erythroplakia}) یا سفتی/زخم در بستر پلاک $\bullet$ (۲) درگیری نواحی پرخطر کف دهان و سطح ونترال/لترال زبان $\bullet$ (۳) لکوپلاکیای وروکوز پرولیفراتیو (\lr{PVL}) چندکانونی مقاوم با ریسک بدخیمی بالای ۶۰٪ $\bullet$ (۴) عدم پسرفت پلاک پس از ۳ تا ۴ هفته از حذف عامل مکانیکی $\bullet$ (۵) لمس لنفادنوپاتی گردنی.}
};

% 10. Refractory / Malignancy Protocols
\node (refractory) [stepbox, line width=0.9pt, text width=15.8cm] at (0, -20.00) {
    \rl{\textbf{پروتکل موارد راجعه، \lr{PVL} و پایش بدخیمی دیسپلازی (\lr{Malignancy Surveillance}):}}\\[3pt]
    \scriptsize \rl{\textbf{سطح ۱ (\lr{PVL} و ضایعات چندکانونی):} نقشه‌برداری بالینی مکرر و بیوپسی‌های اینسیزیونال پیاپی جهت کشف زودرس \lr{Microinvasive OSCC} یا وروکوز کارسینوما \lr{[\textbf{PMID: 40597151}]}.}\\
    \scriptsize \rl{\textbf{سطح ۲ (بررسی حاشیه با فلوئورسنس):} پایش حواشی جراحی با اتوفلوئورسنس نوری (\lr{VELscope}) یا رنگ‌آمیزی تولوییدین بلو جهت تشخیص دیسپلازی ساب‌کلینیکال.}\\
    \scriptsize \rl{\textbf{سطح ۳ (کموتراپی موضعی و فتودینامیک):} کاربرد فتودینامیک‌تراپی (\lr{5-ALA PDT}) \lr{[\textbf{PMID: 42808717}]} یا رتینوئیدهای موضعی؛ در صورت تایید کارسینوم تهاجمی، ارجاع فوری به انکولوژی سر و گردن.}
};

% Connections
\draw [arrow] (n1) -- (n2);
\draw [arrow] (n2) -- (d1);
\draw [arrow] (d1.east) -- node[above, font=\scriptsize\bfseries] {\rl{خاموش / خوش‌خیم}} (r_benign.west);
\draw [arrow] (d1.south) -- node[right=5pt, font=\scriptsize\bfseries] {\rl{فعال / مشکوک}} (n3.north);
\draw [arrow] (n3.south) -- (d2.north);

% Bus-Bar from Decision 2 to 4 Subtype columns
\draw [thick, black] (d2.south) -- (0, -7.50);
\draw [thick, black] (-6.15, -7.50) -- (6.15, -7.50);
\draw [arrow] (-6.15, -7.50) -- (h1.north);
\draw [arrow] (-2.05, -7.50) -- (h2.north);
\draw [arrow] (2.05, -7.50) -- (h3.north);
\draw [arrow, darkred] (6.15, -7.50) -- (h4.north);

% Column Downward Arrows
\draw [arrow] (h1.south) -- (i1.north);
\draw [arrow] (h2.south) -- (i2.north);
\draw [arrow] (h3.south) -- (i3.north);
\draw [arrow, darkred] (h4.south) -- (i4.north);

\draw [arrow] (i1.south) -- (s1.north);
\draw [arrow] (i2.south) -- (s2.north);
\draw [arrow] (i3.south) -- (s3.north);
\draw [arrow, darkred] (i4.south) -- (s4.north);

% Collector Bus into Trigger Box
\draw [thick, black] (s1.south) -- (-6.15, -15.80);
\draw [thick, black] (s2.south) -- (-2.05, -15.80);
\draw [thick, black] (s3.south) -- (2.05, -15.80);
\draw [thick, darkred] (s4.south) -- (6.15, -15.80);
\draw [thick, black] (-6.15, -15.80) -- (6.15, -15.80);
\draw [arrow] (0, -15.80) -- (redflags.north);

% Final link to Refractory Box
\draw [arrow] (redflags.south) -- (refractory.north);
\end{tikzpicture}
\end{latin}
\end{center}

\end{document}
